\chapter{דו-כיווניות בפרסום אקדמי}

פרק זה מוקדש לעקרונות הדו-כיווניות (\textenglish{Bidirectionality}) בפרסום אקדמי
בעברית. הדגשים כוללים שימוש נכון בכיווני טקסט, שילוב מונחים לועזיים בתוך פסקאות
עבריות, והצגת קוד ונוסחאות בסביבה מתאימה.

\section{מנגנון תשומת הלב}

ארכיטקטורת \textenglish{Transformer} מבוססת על מנגנון תשומת הלב
(\textenglish{Attention Mechanism}). מנגנון זה מאפשר לכל אסימון (\textenglish{Token})
ברצף לשים לב לכל שאר האסימונים בו-זמנית, בניגוד לארכיטקטורות רקורנטיות כמו
\textenglish{RNN} ו-\textenglish{LSTM} שעיבדו את הרצף בסדר טורי.

הפונקציה המרכזית היא \textenglish{Scaled Dot-Product Attention}, המחשבת
שקלול של ערכים (\textenglish{Values}) לפי מידת ההתאמה בין שאילתות
(\textenglish{Queries}) ומפתחות (\textenglish{Keys}):

\begin{equation}
    \text{Attention}(Q, K, V) =
    \text{softmax}\!\left(\frac{QK^\top}{\sqrt{d_k}}\right)V
\end{equation}

כאשר \(Q\), \(K\), ו-\(V\) הן מטריצות השאילתה, המפתח והערך בהתאמה,
ו-\(d_k\) הוא מימד המפתח המשמש לנרמול וייצוב הגרדיאנטים.

\section{תשומת לב מרובת ראשים}

\textenglish{Multi-Head Attention} מרחיב את מנגנון תשומת הלב הבסיסי על ידי הפעלת
מספר ראשי תשומת לב במקביל:

\begin{equation}
    \text{MultiHead}(Q, K, V) =
    \text{Concat}(\text{head}_1, \ldots, \text{head}_h) \, W^O
\end{equation}

כאשר כל ראש מוגדר על ידי:

\begin{equation}
    \text{head}_i = \text{Attention}(Q W_i^Q,\; K W_i^K,\; V W_i^V)
\end{equation}

המשמעות היא שהמודל יכול לשים לב בו-זמנית למידע ממיקומים שונים ובמרחבי ייצוג שונים.
מודלים כמו \textenglish{BERT} (\textenglish{Bidirectional Encoder Representations from Transformers})
ו-\textenglish{GPT} (\textenglish{Generative Pre-trained Transformer}) מבוססים על עיקרון זה.

\section{קידוד מיקום}

מכיוון שמנגנון תשומת הלב אינו תלוי-סדר מטבעו, יש להזין מידע מיקומי באופן מפורש.
קידוד המיקום (\textenglish{Positional Encoding}) מוגדר על ידי:

\begin{equation}
    PE_{(pos,\, 2i)}   = \sin\!\left(\frac{pos}{10000^{2i/d_{\text{model}}}}\right)
\end{equation}

\begin{equation}
    PE_{(pos,\, 2i+1)} = \cos\!\left(\frac{pos}{10000^{2i/d_{\text{model}}}}\right)
\end{equation}

כאשר \(pos\) הוא מיקום האסימון ו-\(i\) הוא הממד.

\section{מימוש ב-\textenglish{Python}}

דוגמה למימוש מנגנון תשומת הלב באמצעות \textenglish{PyTorch}:

\begin{LTR}
\begin{verbatim}
import torch
import torch.nn.functional as F
import math

def scaled_dot_product_attention(Q, K, V, mask=None):
    """
    Q, K, V: tensors of shape (batch, heads, seq_len, d_k)
    """
    d_k = Q.size(-1)
    scores = torch.matmul(Q, K.transpose(-2, -1)) / math.sqrt(d_k)
    if mask is not None:
        scores = scores.masked_fill(mask == 0, float('-inf'))
    attn_weights = F.softmax(scores, dim=-1)
    output = torch.matmul(attn_weights, V)
    return output, attn_weights
\end{verbatim}
\end{LTR}

\section{מודלים גדולים של שפה}

מודלים גדולים של שפה (\textenglish{Large Language Models}, קיצור: \textenglish{LLM})
כגון \textenglish{LLaMA}, \textenglish{GPT-4}, ו-\textenglish{Claude}
מאומנים על כמויות עצומות של טקסט. הם מציגים יכולות מרשימות בתחומים רבים כולל
כתיבה יצירתית, פתרון בעיות מתמטיות, ותכנות.

לאימון מודלים אלה נעשה שימוש נרחב במאיצי חומרה (\textenglish{GPU} ו-\textenglish{TPU})
ובמקביליות מבוזרת (\textenglish{Distributed Parallelism}).

\section{דוגמת קוד: ארכיטקטורת \textenglish{Transformer} מלאה}

להלן מימוש מינימלי של שכבת \textenglish{Transformer} בודדת:

\begin{LTR}
\begin{verbatim}
import torch.nn as nn

class TransformerBlock(nn.Module):
    def __init__(self, d_model, num_heads, d_ff, dropout=0.1):
        super().__init__()
        self.attn  = nn.MultiheadAttention(d_model, num_heads,
                                            batch_first=True)
        self.ff    = nn.Sequential(
            nn.Linear(d_model, d_ff),
            nn.ReLU(),
            nn.Linear(d_ff, d_model),
        )
        self.norm1 = nn.LayerNorm(d_model)
        self.norm2 = nn.LayerNorm(d_model)
        self.drop  = nn.Dropout(dropout)

    def forward(self, x, mask=None):
        attn_out, _ = self.attn(x, x, x, attn_mask=mask)
        x = self.norm1(x + self.drop(attn_out))
        x = self.norm2(x + self.drop(self.ff(x)))
        return x
\end{verbatim}
\end{LTR}

\section{סיכום}

פרק זה הדגים את עקרונות הדו-כיווניות בפרסום אקדמי דרך פריזמת ארכיטקטורת
\textenglish{Transformer}. הוצגו נוסחאות מתמטיות, קטעי קוד, ושילוב של מונחים
לועזיים בתוך טקסט עברי — כולם בהתאם לכללי הדו-כיווניות הנכונים.
